\documentclass[12pt, a4paper]{article}

% ─── PAQUETES ────────────────────────────────────────────────────────────────
\usepackage[utf8]{inputenc}
\usepackage[spanish]{babel}
\usepackage[T1]{fontenc}
\usepackage{geometry}
\usepackage{hyperref}
\usepackage{xcolor}
\usepackage{listings}
\usepackage{titlesec}
\usepackage{enumitem}
\usepackage{fancyhdr}
\usepackage{graphicx}
\usepackage{booktabs}
\usepackage{array}
\usepackage{longtable}
\usepackage{mdframed}
\usepackage{tcolorbox}
\usepackage{fontawesome5}
\usepackage{parskip}
\usepackage{microtype}
\usepackage{caption}
\usepackage{float}
\usepackage{tikz}

% ─── GEOMETRÍA ───────────────────────────────────────────────────────────────
\geometry{
    top=2.5cm,
    bottom=2.5cm,
    left=2.8cm,
    right=2.8cm
}

% ─── COLORES DEL PROYECTO ────────────────────────────────────────────────────
\definecolor{cpAccent}{HTML}{6C63FF}       % Morado principal
\definecolor{cpFondo}{HTML}{0F1117}        % Fondo oscuro
\definecolor{cpTarjeta}{HTML}{1A1D27}      % Fondo tarjeta
\definecolor{cpTextoClaro}{HTML}{8B8FA3}   % Texto secundario
\definecolor{cpVerde}{HTML}{4CAF50}        % Verde
\definecolor{cpAzul}{HTML}{2196F3}         % Azul
\definecolor{cpNaranja}{HTML}{FF9800}      % Naranja
\definecolor{cpKeyword}{HTML}{C792EA}      % Palabras clave (morado claro)
\definecolor{cpString}{HTML}{C3E88D}       % Strings (verde claro)
\definecolor{cpComment}{HTML}{546E7A}      % Comentarios (gris)
\definecolor{cpFunction}{HTML}{82AAFF}     % Funciones (azul)
\definecolor{cpNumber}{HTML}{F78C6C}       % Números (naranja)
\definecolor{cpBuiltin}{HTML}{FFCB6B}      % Builtins (amarillo)
\definecolor{cpDecorator}{HTML}{FF5370}    % Decoradores (rojo)
\definecolor{cpWarn}{HTML}{FFF3CD}         % Fondo advertencia
\definecolor{cpInfo}{HTML}{D1ECF1}         % Fondo info
\definecolor{cpGray}{HTML}{F8F9FA}         % Fondo gris claro

% ─── CAJAS DE COLORES ────────────────────────────────────────────────────────
\tcbuselibrary{skins, breakable, listings}

\newtcolorbox{cajaInfo}{
    colback=cpInfo,
    colframe=cpAzul,
    fonttitle=\bfseries,
    arc=4pt,
    left=8pt, right=8pt, top=6pt, bottom=6pt
}

\newtcolorbox{cajaAviso}{
    colback=cpWarn,
    colframe=cpNaranja,
    fonttitle=\bfseries,
    arc=4pt,
    left=8pt, right=8pt, top=6pt, bottom=6pt
}

\newtcolorbox{cajaCodigo}{
    colback=cpGray,
    colframe=black!15,
    arc=4pt,
    left=8pt, right=8pt, top=6pt, bottom=6pt,
    fontupper=\ttfamily\small
}

% ─── CONFIGURACIÓN DE LISTINGS (código fuente) ───────────────────────────────
\lstdefinestyle{htmlstyle}{
    language=HTML,
    basicstyle=\ttfamily\footnotesize,
    keywordstyle=\color{cpKeyword}\bfseries,
    stringstyle=\color{cpString},
    commentstyle=\color{cpComment}\itshape,
    backgroundcolor=\color{cpGray},
    frame=single,
    framesep=4pt,
    rulecolor=\color{black!20},
    breaklines=true,
    breakatwhitespace=true,
    showstringspaces=false,
    tabsize=4,
    numbers=left,
    numberstyle=\tiny\color{cpTextoClaro},
    numbersep=8pt,
    xleftmargin=15pt,
    captionpos=b
}

\lstdefinestyle{pythonstyle}{
    language=Python,
    basicstyle=\ttfamily\footnotesize,
    keywordstyle=\color{cpKeyword}\bfseries,
    stringstyle=\color{cpString},
    commentstyle=\color{cpComment}\itshape,
    backgroundcolor=\color{cpGray},
    frame=single,
    framesep=4pt,
    rulecolor=\color{black!20},
    breaklines=true,
    breakatwhitespace=true,
    showstringspaces=false,
    tabsize=4,
    numbers=left,
    numberstyle=\tiny\color{cpTextoClaro},
    numbersep=8pt,
    xleftmargin=15pt,
    captionpos=b
}

\lstdefinestyle{cssInline}{
    language=HTML,
    basicstyle=\ttfamily\footnotesize,
    backgroundcolor=\color{cpGray},
    frame=single,
    framesep=4pt,
    rulecolor=\color{black!20},
    breaklines=true,
    showstringspaces=false,
    tabsize=2,
    xleftmargin=10pt,
    captionpos=b
}

% ─── ENCABEZADOS Y PIE ───────────────────────────────────────────────────────
\pagestyle{fancy}
\fancyhf{}
\fancyhead[L]{\textcolor{cpAccent}{\textbf{ConversionPulse}} \ -- \ Documentación Técnica}
\fancyhead[R]{\textcolor{cpTextoClaro}{\small v1.0 · 2025}}
\fancyfoot[C]{\textcolor{cpTextoClaro}{\small \thepage}}
\renewcommand{\headrulewidth}{0.4pt}
\renewcommand{\footrulewidth}{0pt}
\renewcommand{\headrule}{\hbox to\headwidth{\color{cpAccent}\leaders\hrule height\headrulewidth\hfill}}

% ─── ESTILO DE SECCIONES ─────────────────────────────────────────────────────
\titleformat{\section}
    {\large\bfseries\color{cpAccent}}
    {\thesection.}{0.6em}{}
    [\vspace{-4pt}\textcolor{cpAccent}{\rule{\linewidth}{0.6pt}}]

\titleformat{\subsection}
    {\normalsize\bfseries\color{black!80}}
    {\thesubsection.}{0.5em}{}

\titleformat{\subsubsection}
    {\small\bfseries\color{black!70}}
    {\thesubsubsection.}{0.5em}{}

% ─── HIPERLINKS ──────────────────────────────────────────────────────────────
\hypersetup{
    colorlinks=true,
    linkcolor=cpAccent,
    urlcolor=cpAzul,
    citecolor=cpVerde,
    pdftitle={ConversionPulse - Documentación Técnica},
    pdfauthor={ConversionPulse Team},
    pdfsubject={Repositorio interno de código HTML/CSS},
}

% ─── COMANDOS ÚTILES ─────────────────────────────────────────────────────────
\newcommand{\archivo}[1]{\texttt{\textcolor{cpAccent}{#1}}}
\newcommand{\clase}[1]{\texttt{\textcolor{cpAzul}{.#1}}}
\newcommand{\etiquetaHtml}[1]{\texttt{\textcolor{cpKeyword}{<#1>}}}
\newcommand{\variable}[1]{\texttt{\textcolor{cpNumber}{--#1}}}
\newcommand{\codigo}[1]{\texttt{\small #1}}
\newcommand{\badge}[2]{\fcolorbox{#1!40}{#1!15}{\textcolor{#1}{\small\textbf{#2}}}}


% ═════════════════════════════════════════════════════════════════════════════
%  DOCUMENTO
% ═════════════════════════════════════════════════════════════════════════════
\begin{document}

% ─── PORTADA ─────────────────────────────────────────────────────────────────
\begin{titlepage}
    \centering
    \vspace*{2cm}

    {\Huge\bfseries\color{cpAccent} ConversionPulse}\\[0.5em]
    {\Large\color{black!70} Repositorio Interno de Código}

    \vspace{1.5cm}
    \noindent\textcolor{cpAccent}{\rule{10cm}{2pt}}
    \vspace{1.5cm}

    {\LARGE\bfseries Documentación Técnica Completa}\\[0.8em]
    {\large Guía de Edición, Mantenimiento y Expansión}

    \vspace{2cm}
    \begin{tabular}{ll}
        \textbf{Versión:}    & 1.0 \\[4pt]
        \textbf{Fecha:}      & Mayo 2025 \\[4pt]
        \textbf{Tecnologías:} & HTML5 · CSS3 · JavaScript (vanilla) \\[4pt]
        \textbf{Códigos:}    & 16 entradas en 4 categorías \\[4pt]
        \textbf{Propósito:}  & Uso interno — Repositorio de demos \\
    \end{tabular}

    \vspace{2cm}
    \begin{cajaAviso}
        \textbf{Importante:} Este documento es la guía de referencia para cualquier
        persona que necesite agregar, editar o modificar el repositorio.
        Léelo completo antes de hacer cualquier cambio.
    \end{cajaAviso}

    \vfill
    {\small\color{cpTextoClaro} ConversionPulse · Repositorio Interno · 2025}
\end{titlepage}

% ─── TABLA DE CONTENIDOS ─────────────────────────────────────────────────────
\tableofcontents
\newpage


% ═════════════════════════════════════════════════════════════════════════════
\section{Visión General del Proyecto}
% ═════════════════════════════════════════════════════════════════════════════

\subsection{¿Qué es ConversionPulse?}

ConversionPulse es un \textbf{repositorio interno de código} presentado como un
sitio web estático (sin servidor, sin base de datos, sin frameworks). Su propósito
es centralizar, documentar y hacer accesibles snippets, scripts, notebooks y
proyectos Python organizados por caso de uso, para que los equipos internos
puedan consultarlos y reutilizarlos.

El sitio funciona directamente en el navegador abriendo el archivo
\archivo{index.html}. No requiere instalación, dependencias ni conexión a internet.

\subsection{Características principales}

\begin{itemize}[leftmargin=1.5em]
    \item \textbf{16 códigos} distribuidos en 4 categorías de uso.
    \item \textbf{Resaltado de sintaxis} manual con clases CSS (no requiere librerías externas).
    \item \textbf{Botón de copiar} código con un clic.
    \item \textbf{Soporte de video}: YouTube (iframe) o archivos \codigo{.mp4} locales.
    \item \textbf{Diseño oscuro} consistente con variables CSS modificables.
    \item \textbf{Buscador} en el sidebar (funcional con JavaScript).
    \item \textbf{Responsive} en escritorio; el layout es sidebar fijo + contenido central.
\end{itemize}

\subsection{Tecnologías utilizadas}

\begin{center}
\begin{tabular}{lll}
\toprule
\textbf{Tecnología} & \textbf{Versión} & \textbf{Uso} \\
\midrule
HTML5  & --    & Estructura de todas las páginas \\
CSS3   & --    & Estilos centralizados en \archivo{styles.css} \\
JavaScript & Vanilla (sin librerías) & Botón copiar, buscador del sidebar \\
\bottomrule
\end{tabular}
\end{center}

\begin{cajaInfo}
    \textbf{Sin dependencias externas:} No se usa React, Vue, Bootstrap,
    jQuery ni ninguna librería. Todo corre directamente en el navegador
    sin necesidad de \texttt{npm install} ni servidores.
\end{cajaInfo}


% ═════════════════════════════════════════════════════════════════════════════
\section{Estructura de Archivos}
% ═════════════════════════════════════════════════════════════════════════════

\subsection{Árbol de directorios}

\begin{cajaCodigo}
repo-demos/\\
\quad├── index.html             \textit{← Dashboard principal}\\
\quad├── styles.css             \textit{← Todos los estilos del sitio}\\
\quad├── PLANTILLA.html         \textit{← Plantilla para nuevos códigos}\\
\quad├── documentacion.tex      \textit{← Este documento}\\
\quad│\\
\quad├── categorias/            \textit{← Páginas de cada categoría}\\
\quad│   ├── analytics.html\\
\quad│   ├── financial-crimes.html\\
\quad│   ├── gobierno.html\\
\quad│   └── retail.html\\
\quad│\\
\quad└── codigos/               \textit{← Una página por código}\\
\quad    ├── ana-001.html        \textit{Pipeline de ETL Automatizado}\\
\quad    ├── ana-002.html        \textit{A/B Testing Framework}\\
\quad    ├── ana-003.html        \textit{Dashboard de Métricas en Tiempo Real}\\
\quad    ├── ana-004.html        \textit{Modelo de Atribución Multi-Touch}\\
\quad    ├── fc-001.html         \textit{Detección de Transacciones Anómalas}\\
\quad    ├── fc-002.html         \textit{Red de Lavado de Dinero - Graph Analysis}\\
\quad    ├── fc-003.html         \textit{Score de Riesgo - Credit Card Fraud}\\
\quad    ├── fc-004.html         \textit{KYC Automation Pipeline}\\
\quad    ├── gob-001.html        \textit{Análisis de Presupuesto Público}\\
\quad    ├── gob-002.html        \textit{Predicción de Demanda de Servicios}\\
\quad    ├── gob-003.html        \textit{Dashboard de Indicadores Socioeconómicos}\\
\quad    ├── gob-004.html        \textit{Clasificación de Documentos Oficiales}\\
\quad    ├── ret-003.html        \textit{Motor de Recomendaciones Colaborativo}\\
\quad    ├── ret-004.html        \textit{Predicción de Demanda de Inventario}\\
\quad    ├── ret-005.html        \textit{Segmentación de Clientes (RFM)}\\
\quad    └── ret-006.html        \textit{Market Basket Analysis}
\end{cajaCodigo}

\subsection{Convención de nombres de archivos}

Los archivos de la carpeta \archivo{codigos/} siguen el patrón:

\begin{cajaCodigo}
\textbf{[prefijo]-[número].html}\\[4pt]
Ejemplos: \texttt{ana-001.html}, \texttt{fc-003.html}, \texttt{gob-002.html}
\end{cajaCodigo}

Los prefijos disponibles son:

\begin{center}
\begin{tabular}{cll}
\toprule
\textbf{Prefijo} & \textbf{Categoría} & \textbf{Icono} \\
\midrule
\texttt{ana} & Analytics           & 📈 \\
\texttt{fc}  & Financial Crimes    & 🔒 \\
\texttt{gob} & Gobierno            & 🏛️ \\
\texttt{ret} & Retail              & 🛒 \\
\bottomrule
\end{tabular}
\end{center}

\begin{cajaAviso}
    \textbf{Nota sobre Retail:} Los archivos de Retail inician en
    \archivo{ret-003.html}. Los códigos \archivo{ret-001.html} y
    \archivo{ret-002.html} no existen aún — están pendientes de crear.
    Al agregar nuevos Retail, continúa desde \archivo{ret-007.html}.
\end{cajaAviso}


% ═════════════════════════════════════════════════════════════════════════════
\section{Catálogo de Códigos}
% ═════════════════════════════════════════════════════════════════════════════

\subsection{Analytics (ANA)}

\begin{longtable}{lllp{5.5cm}}
\toprule
\textbf{ID} & \textbf{Tipo} & \textbf{Autor} & \textbf{Título / Tags} \\
\midrule
\endhead
\texttt{ana-001} & Proyecto  & Data Engineering       & \textit{Pipeline de ETL Automatizado}\\
                 &           &                        & \texttt{\#etl \#pipeline \#airflow} \\[6pt]
\texttt{ana-002} & Notebook  & Experimentation Team   & \textit{A/B Testing Framework}\\
                 &           &                        & \texttt{\#ab-testing \#statistics \#hypothesis-testing} \\[6pt]
\texttt{ana-003} & Script    & Platform Team          & \textit{Dashboard de Métricas en Tiempo Real}\\
                 &           &                        & \texttt{\#dashboard \#streaming \#kafka} \\[6pt]
\texttt{ana-004} & Notebook  & Marketing Analytics    & \textit{Modelo de Atribución Multi-Touch}\\
                 &           &                        & \texttt{\#attribution \#marketing \#multi-touch} \\
\bottomrule
\end{longtable}

\subsection{Financial Crimes (FC)}

\begin{longtable}{lllp{5.5cm}}
\toprule
\textbf{ID} & \textbf{Tipo} & \textbf{Autor} & \textbf{Título / Tags} \\
\midrule
\endhead
\texttt{fc-001} & Notebook  & Data Science Team         & \textit{Detección de Transacciones Anómalas}\\
                &           &                           & \texttt{\#fraud-detection \#anomaly-detection \#isolation-forest} \\[6pt]
\texttt{fc-002} & Script    & Compliance Team           & \textit{Red de Lavado de Dinero --- Graph Analysis}\\
                &           &                           & \texttt{\#money-laundering \#graph-analysis \#networkx} \\[6pt]
\texttt{fc-003} & Proyecto  & ML Engineering            & \textit{Score de Riesgo --- Credit Card Fraud}\\
                &           &                           & \texttt{\#credit-card \#risk-scoring \#real-time} \\[6pt]
\texttt{fc-004} & Script    & Compliance Engineering    & \textit{KYC Automation Pipeline}\\
                &           &                           & \texttt{\#kyc \#automation \#document-verification} \\
\bottomrule
\end{longtable}

\subsection{Gobierno (GOB)}

\begin{longtable}{lllp{5.5cm}}
\toprule
\textbf{ID} & \textbf{Tipo} & \textbf{Autor} & \textbf{Título / Tags} \\
\midrule
\endhead
\texttt{gob-001} & Notebook & Policy Analytics      & \textit{Análisis de Presupuesto Público}\\
                 &          &                       & \texttt{\#budget-analysis \#visualization \#pandas} \\[6pt]
\texttt{gob-002} & Proyecto & Government Solutions  & \textit{Predicción de Demanda de Servicios}\\
                 &          &                       & \texttt{\#demand-forecasting \#time-series \#arima} \\[6pt]
\texttt{gob-003} & Script   & Data Governance       & \textit{Dashboard de Indicadores Socioeconómicos}\\
                 &          &                       & \texttt{\#dashboard \#indicators \#plotly} \\[6pt]
\texttt{gob-004} & Notebook & NLP Team              & \textit{Clasificación de Documentos Oficiales}\\
                 &          &                       & \texttt{\#nlp \#text-classification \#transformers} \\
\bottomrule
\end{longtable}

\subsection{Retail (RET)}

\begin{longtable}{lllp{5.5cm}}
\toprule
\textbf{ID} & \textbf{Tipo} & \textbf{Autor} & \textbf{Título / Tags} \\
\midrule
\endhead
\texttt{ret-003} & Proyecto & Personalization Team  & \textit{Motor de Recomendaciones Colaborativo}\\
                 &          &                       & \texttt{\#recommendation-engine \#collaborative-filtering} \\[6pt]
\texttt{ret-004} & Notebook & Supply Chain Team     & \textit{Predicción de Demanda de Inventario}\\
                 &          &                       & \texttt{\#demand-forecast \#inventory \#prophet} \\[6pt]
\texttt{ret-005} & Script   & CRM Analytics         & \textit{Segmentación de Clientes (RFM)}\\
                 &          &                       & \texttt{\#rfm \#segmentation \#clustering} \\[6pt]
\texttt{ret-006} & Notebook & Retail Analytics      & \textit{Market Basket Analysis}\\
                 &          &                       & \texttt{\#market-basket \#apriori \#association-rules} \\
\bottomrule
\end{longtable}


% ═════════════════════════════════════════════════════════════════════════════
\section{Anatomía de una Página de Código}
% ═════════════════════════════════════════════════════════════════════════════

Cada archivo en \archivo{codigos/} tiene la misma estructura HTML.
Entenderla es fundamental para editar o crear entradas.

\subsection{Estructura general del HTML}

\begin{lstlisting}[style=htmlstyle, caption={Esqueleto de una página de código}]
<!DOCTYPE html>
<html lang="es">
<head>
    <meta charset="UTF-8">
    <title>ConversionPulse - NOMBRE DEL PROYECTO</title>
    <link rel="stylesheet" href="../styles.css">   <!-- Nunca cambiar la ruta -->
</head>
<body>

    <nav class="sidebar"> ... </nav>      <!-- Sidebar: NO tocar -->

    <div class="zona-principal">
        <div class="header"> ... </div>   <!-- Header: NO tocar -->

        <div class="contenido">
            <a class="boton-volver">...</a>        <!-- Volver al index -->

            <div class="tarjeta">
                <!-- [1] TITULO -->
                <h1 class="tarjeta-titulo">NOMBRE DEL PROYECTO</h1>

                <!-- [2] LENGUAJE Y TIPO -->
                <span class="etiqueta-lenguaje">python</span>
                <span>Proyecto</span>              <!-- o Notebook, Script -->

                <!-- [3] DESCRIPCION -->
                <p class="tarjeta-descripcion">...</p>

                <!-- [4] ETIQUETAS / TAGS -->
                <div class="etiquetas">
                    <span class="etiqueta">#tag1</span>
                </div>

                <!-- [5] AUTOR Y FECHA -->
                <div class="info-meta">
                    <span><strong>Autor:</strong> Equipo</span>
                    <span><strong>Fecha:</strong> 2025-01-01</span>
                </div>

                <!-- [6] INSTRUCCIONES -->
                <div class="instrucciones">
                    <h2>Instrucciones de Uso</h2>
                    <p>...</p>
                </div>

                <!-- [7] BLOQUE DE CODIGO -->
                <div class="bloque-codigo">
                    <div class="barra-archivo">
                        <span>nombre_archivo.py</span>
                        <span onclick="copiarCodigo(this)">Copiar</span>
                    </div>
                    <pre><code>...codigo aqui...</code></pre>
                </div>

                <!-- [8] VIDEO (opcional) -->
                <div class="seccion-video"> ... </div>
            </div>
        </div>

        <div class="footer"> ... </div>
    </div>

    <script>
        function copiarCodigo(boton) { ... }    <!-- No tocar -->
    </script>
</body>
</html>
\end{lstlisting}

\subsection{Elementos obligatorios vs opcionales}

\begin{center}
\begin{tabular}{lcc}
\toprule
\textbf{Elemento} & \textbf{¿Obligatorio?} & \textbf{Notas} \\
\midrule
\etiquetaHtml{title} en \etiquetaHtml{head}     & Sí & \texttt{ConversionPulse - NOMBRE} \\
\clase{tarjeta-titulo}                           & Sí & Título visible grande \\
\clase{etiqueta-lenguaje}                        & Sí & \texttt{python}, \texttt{sql}, etc. \\
Tipo (Proyecto / Notebook / Script)              & Sí & Texto simple entre spans \\
\clase{tarjeta-descripcion}                      & Sí & 1--2 oraciones \\
\clase{etiquetas}                                & Sí & Mínimo 1 etiqueta \\
\clase{info-meta} (autor + fecha)                & Sí & Formato \texttt{YYYY-MM-DD} \\
\clase{instrucciones}                            & Sí & Qué recibe, qué retorna \\
\clase{bloque-codigo}                            & Sí & Al menos uno \\
\clase{seccion-video}                            & No & Puede quedar el placeholder \\
Múltiples bloques de código                      & No & Copiar el bloque y descomentar \\
\bottomrule
\end{tabular}
\end{center}

\subsection{El sidebar y qué clase ``activa'' usar}

El sidebar aparece en todas las páginas. El único cambio que varía entre
archivos es cuál categoría tiene la clase \codigo{activa}:

\begin{lstlisting}[style=htmlstyle, caption={Activar la categoría correcta en el sidebar}]
<!-- En una pagina de Analytics: -->
<a href="../categorias/analytics.html" class="sidebar-categoria activa">

<!-- En una pagina de Financial Crimes: -->
<a href="../categorias/financial-crimes.html" class="sidebar-categoria activa">

<!-- En una pagina de Gobierno: -->
<a href="../categorias/gobierno.html" class="sidebar-categoria activa">

<!-- En una pagina de Retail: -->
<a href="../categorias/retail.html" class="sidebar-categoria activa">
\end{lstlisting}

\begin{cajaAviso}
    \textbf{Rutas relativas:} Los archivos en \archivo{codigos/} usan
    \texttt{../} para subir un nivel. Nunca uses rutas absolutas como
    \texttt{/categorias/...} ya que romperán al abrir el sitio localmente.
\end{cajaAviso}


% ═════════════════════════════════════════════════════════════════════════════
\section{Sistema de Resaltado de Sintaxis}
% ═════════════════════════════════════════════════════════════════════════════

El sitio no usa ninguna librería de resaltado de sintaxis (como Prism.js o
Highlight.js). En su lugar, el código se colorea manualmente envolviendo
palabras en etiquetas \etiquetaHtml{span} con clases CSS específicas.

\subsection{Tabla de clases de color}

\begin{center}
\begin{tabular}{llll}
\toprule
\textbf{Clase CSS} & \textbf{Color} & \textbf{Hex} & \textbf{Usar para...} \\
\midrule
\clase{keyword}   & Morado claro & \texttt{\#C792EA} & \texttt{import, class, def, if, return, from, as, not, in, for} \\
\clase{string}    & Verde claro  & \texttt{\#C3E88D} & Textos entre comillas simples o dobles \\
\clase{comment}   & Gris         & \texttt{\#546E7A} & Líneas que empiezan con \texttt{\#} \\
\clase{function}  & Azul         & \texttt{\#82AAFF} & Nombres de funciones y clases \\
\clase{number}    & Naranja      & \texttt{\#F78C6C} & Valores numéricos (\texttt{42}, \texttt{3.14}, \texttt{True}, \texttt{False}, \texttt{None}) \\
\clase{builtin}   & Amarillo     & \texttt{\#FFCB6B} & Funciones built-in: \texttt{print}, \texttt{len}, \texttt{range}, \texttt{float}, \texttt{int} \\
\clase{decorator} & Rojo         & \texttt{\#FF5370} & Decoradores: \texttt{@staticmethod}, \texttt{@property} \\
\clase{param}     & Rosa         & \texttt{\#F07178} & Parámetros de funciones \\
\bottomrule
\end{tabular}
\end{center}

\subsection{Ejemplo de código coloreado}

Así es como se escribe código coloreado dentro del HTML:

\begin{lstlisting}[style=htmlstyle, caption={Ejemplo: función Python con colores}]
<pre><code><span class="keyword">import</span> pandas <span class="keyword">as</span> pd

<span class="keyword">class</span> <span class="function">MiClase</span>:
    <span class="string">"""Docstring de la clase"""</span>

    <span class="keyword">def</span> <span class="function">mi_metodo</span>(self, valor=<span class="number">10</span>):
        <span class="comment"># Este es un comentario</span>
        resultado = valor * <span class="number">2</span>
        <span class="builtin">print</span>(<span class="string">"Resultado:"</span>, resultado)
        <span class="keyword">return</span> resultado

obj = <span class="function">MiClase</span>()
obj.mi_metodo(<span class="number">5</span>)</code></pre>
\end{lstlisting}

\subsection{Estrategia para colorear código nuevo}

Sigue este orden para colorear código de forma eficiente:

\begin{enumerate}[leftmargin=1.5em]
    \item Pega el código limpio dentro de \etiquetaHtml{pre}\etiquetaHtml{code}.
    \item Reemplaza los \textbf{comentarios} (líneas con \texttt{\#}) con \clase{comment}.
    \item Reemplaza las \textbf{palabras clave} (\texttt{import}, \texttt{def}, \texttt{class}, etc.) con \clase{keyword}.
    \item Reemplaza los \textbf{strings} (texto entre comillas) con \clase{string}.
    \item Reemplaza los \textbf{nombres de funciones} con \clase{function}.
    \item Reemplaza los \textbf{números} con \clase{number}.
    \item Reemplaza las \textbf{funciones built-in} con \clase{builtin}.
\end{enumerate}

\begin{cajaInfo}
    \textbf{¿El código no necesita colores?} Es válido pegar el código
    directamente en texto plano dentro de \etiquetaHtml{pre}\etiquetaHtml{code}
    sin ningún \texttt{span}. Se verá en gris claro sobre fondo oscuro,
    lo cual es perfectamente legible.
\end{cajaInfo}

\subsection{Caracteres especiales en HTML}

Si el código contiene los caracteres \texttt{<}, \texttt{>} o \texttt{\&},
deben escaparse para que el navegador no los interprete como HTML:

\begin{center}
\begin{tabular}{lll}
\toprule
\textbf{Carácter} & \textbf{Código HTML} & \textbf{Ejemplo} \\
\midrule
\texttt{<}  & \texttt{\&lt;}   & \texttt{a \&lt; b} → muestra \texttt{a < b} \\
\texttt{>}  & \texttt{\&gt;}   & \texttt{a \&gt; b} → muestra \texttt{a > b} \\
\texttt{\&} & \texttt{\&amp;}  & \texttt{x \&amp;\&amp; y} → muestra \texttt{x \&\& y} \\
\bottomrule
\end{tabular}
\end{center}


% ═════════════════════════════════════════════════════════════════════════════
\section{Agregar un Nuevo Código: Paso a Paso}
% ═════════════════════════════════════════════════════════════════════════════

\subsection{Paso 1: Copiar la plantilla}

Copia el archivo \archivo{PLANTILLA.html} y renómbralo según la convención:

\begin{cajaCodigo}
\textbf{Formato:} \texttt{[prefijo]-[número].html}\\[4pt]
Ejemplo: para el quinto Analytics → \texttt{ana-005.html}\\
Guárdalo dentro de la carpeta \texttt{codigos/}
\end{cajaCodigo}

\subsection{Paso 2: Actualizar el \texttt{<title>}}

En la línea del \etiquetaHtml{head}, cambia el título:

\begin{lstlisting}[style=htmlstyle]
<title>ConversionPulse - NOMBRE DE TU PROYECTO</title>
\end{lstlisting}

\subsection{Paso 3: Activar la categoría correcta en el sidebar}

Busca la sección del sidebar y agrega \texttt{activa} al enlace de la categoría
correspondiente (ver sección 4.3 de este documento).

\subsection{Paso 4: Rellenar la tarjeta}

Edita cada sección marcada con \texttt{CAMBIAR} en la plantilla:

\begin{enumerate}[leftmargin=1.5em]
    \item \textbf{\clase{tarjeta-titulo}:} Nombre del proyecto (ej: ``Análisis de Churn'').
    \item \textbf{\clase{etiqueta-lenguaje}:} Lenguaje de programación (ej: \texttt{python}, \texttt{sql}, \texttt{r}).
    \item \textbf{Tipo:} Cambia el texto \texttt{Proyecto} por \texttt{Notebook} o \texttt{Script} según corresponda.
    \item \textbf{\clase{tarjeta-descripcion}:} 1--2 oraciones que expliquen qué hace el código.
    \item \textbf{\clase{etiquetas}:} 2--4 tags relevantes. Usa kebab-case: \texttt{\#machine-learning}.
    \item \textbf{\clase{info-meta}:} Autor (nombre o equipo) y fecha en formato \texttt{YYYY-MM-DD}.
    \item \textbf{\clase{instrucciones}:} Explica qué datos recibe, qué retorna, qué dependencias necesita.
\end{enumerate}

\subsection{Paso 5: Agregar el código}

Dentro del \clase{bloque-codigo}:

\begin{enumerate}[leftmargin=1.5em]
    \item Cambia el nombre del archivo en la \clase{barra-archivo}: \texttt{nombre\_archivo.py}
    \item Pega tu código dentro de \etiquetaHtml{pre}\etiquetaHtml{code}...\etiquetaHtml{/code}\etiquetaHtml{/pre}
    \item Opcionalmente, coloréalo con las clases de la sección 5.
\end{enumerate}

\textbf{¿Necesitas más de un bloque de código?} La plantilla incluye un segundo
bloque comentado. Para activarlo, elimina los marcadores \texttt{<!--} y \texttt{-->}
que lo rodean:

\begin{lstlisting}[style=htmlstyle, caption={Activar un segundo bloque de código}]
<!-- Antes (comentado = inactivo): -->
<!--
<div class="bloque-codigo">
    ...
</div>
-->

<!-- Despues (sin comentarios = activo): -->
<div class="bloque-codigo">
    ...
</div>
\end{lstlisting}

\subsection{Paso 6: Agregar video (opcional)}

Dentro de \clase{seccion-video}, reemplaza el \clase{video-placeholder}
con una de estas dos opciones:

\textbf{Opción A — Video de YouTube:}
\begin{lstlisting}[style=htmlstyle]
<div class="video-wrapper">
    <iframe
        src="https://www.youtube.com/embed/TU_VIDEO_ID"
        allowfullscreen>
    </iframe>
</div>
\end{lstlisting}

\begin{cajaInfo}
    Para obtener el ID del video de YouTube, toma la URL del video:
    \texttt{https://www.youtube.com/watch?v=\textbf{dQw4w9WgXcQ}}
    El ID es la parte después de \texttt{v=}: \texttt{dQw4w9WgXcQ}
\end{cajaInfo}

\textbf{Opción B — Archivo de video local:}
\begin{lstlisting}[style=htmlstyle]
<div class="video-wrapper">
    <video controls>
        <source src="../videos/mi_video.mp4" type="video/mp4">
        Tu navegador no soporta video HTML5.
    </video>
</div>
\end{lstlisting}

\begin{cajaAviso}
    Si usas videos locales, crea una carpeta \archivo{videos/} en la raíz
    del proyecto y coloca los archivos \texttt{.mp4} ahí. El tamaño
    recomendado es menor a 50 MB por video.
\end{cajaAviso}

\subsection{Paso 7: Registrar la entrada en la página de categoría}

Ve al archivo correspondiente en \archivo{categorias/} (ej: \archivo{analytics.html})
y agrega una tarjeta en el \clase{grid-proyectos}:

\begin{lstlisting}[style=htmlstyle, caption={Agregar tarjeta en página de categoría}]
<a href="../codigos/ana-005.html" class="proyecto-card">
    <div class="card-top">
        <span class="card-titulo">Nombre del Proyecto</span>
        <span class="badge-tipo badge-notebook">Notebook</span>
        <!-- Tipos de badge: badge-notebook | badge-script | badge-proyecto -->
    </div>
    <div class="card-desc">
        Descripcion breve del proyecto (igual que la tarjeta-descripcion).
    </div>
    <div class="card-tags">
        <span class="card-tag">tag1</span>
        <span class="card-tag">tag2</span>
        <span class="card-lang">
            <span class="etiqueta-lenguaje">python</span>
        </span>
    </div>
    <div class="card-footer">
        <span>Nombre del Equipo</span>
        <span>2025-01-01</span>
    </div>
</a>
\end{lstlisting}

\subsection{Paso 8: Actualizar los contadores}

Hay \textbf{tres lugares} donde se muestra el número total de códigos.
Actualízalos cuando agregues o elimines entradas:

\begin{center}
\begin{tabular}{lll}
\toprule
\textbf{Archivo} & \textbf{Elemento} & \textbf{Qué actualizar} \\
\midrule
\archivo{index.html}       & \clase{stat-numero} (primero)       & Total de códigos globales \\
\archivo{index.html}       & \clase{sidebar-footer} → \clase{detalle} & Total en el sidebar \\
\archivo{index.html}       & \clase{cat-count} de la categoría   & Cantidad por categoría \\
\archivo{categorias/*.html} & \clase{categoria-contador}          & ``X códigos disponibles'' \\
Todos los \archivo{codigos/*.html} & \clase{sidebar-footer} → \clase{detalle} & ``16 códigos disponibles'' \\
\bottomrule
\end{tabular}
\end{center}

\begin{cajaAviso}
    \textbf{Detalle importante:} El texto ``16 códigos disponibles'' aparece
    en el \texttt{sidebar-footer} de \textbf{cada archivo} de código. Si
    agregas un nuevo código y el total pasa a 17, debes editar ese texto
    en todos los archivos de la carpeta \archivo{codigos/}.
\end{cajaAviso}

\subsection{Paso 9: Agregar a la lista de ``Códigos Recientes'' (opcional)}

En \archivo{index.html}, la sección \clase{dash-recientes} muestra los 5 códigos
más recientes. Si quieres que tu nuevo código aparezca ahí:

\begin{lstlisting}[style=htmlstyle, caption={Agregar entrada a Códigos Recientes en index.html}]
<li>
    <a href="codigos/ana-005.html">
        <span class="rec-icono">📊</span>   <!-- Elige un emoji relevante -->
        <div class="rec-info">
            <div class="rec-nombre">Nombre del Proyecto</div>
            <div class="rec-desc">Descripcion breve que se corta con puntos suspensivos...</div>
        </div>
        <div class="rec-meta">
            <span class="etiqueta-lenguaje">python</span>
            <span class="rec-fecha">2025-01-01</span>
        </div>
    </a>
</li>
\end{lstlisting}


% ═════════════════════════════════════════════════════════════════════════════
\section{Crear una Nueva Categoría}
% ═════════════════════════════════════════════════════════════════════════════

Si necesitas agregar una quinta categoría (ej: ``Machine Learning''), sigue
estos pasos:

\subsection{Paso 1: Crear el archivo de categoría}

Copia cualquier archivo existente en \archivo{categorias/} (ej: \archivo{analytics.html})
y renómbralo (ej: \archivo{machine-learning.html}). Luego edita:

\begin{itemize}[leftmargin=1.5em]
    \item El \etiquetaHtml{title}: \texttt{ConversionPulse - Machine Learning}
    \item El \clase{categoria-icono}: cambia el emoji
    \item El \etiquetaHtml{h1}: nombre de la categoría
    \item La \clase{categoria-descripcion}: descripción de la categoría
    \item Vacía el \clase{grid-proyectos}: elimina las tarjetas existentes
    \item Actualiza el link activo en el sidebar
\end{itemize}

\subsection{Paso 2: Agregar la categoría al sidebar (en todos los archivos)}

En \textbf{todos} los archivos HTML del proyecto, agrega el nuevo enlace al sidebar:

\begin{lstlisting}[style=htmlstyle]
<a href="../categorias/machine-learning.html" class="sidebar-categoria">
    <span class="icono">🤖</span> Machine Learning
</a>
\end{lstlisting}

\begin{cajaAviso}
    Recuerda que en \archivo{index.html} la ruta es diferente (sin \texttt{../}):
    \texttt{href="categorias/machine-learning.html"}
\end{cajaAviso}

\subsection{Paso 3: Agregar la tarjeta de categoría en el dashboard}

En \archivo{index.html}, dentro de \clase{dash-categorias}:

\begin{lstlisting}[style=htmlstyle]
<a href="categorias/machine-learning.html" class="dash-cat-card">
    <span class="cat-icono">🤖</span>
    <span class="cat-nombre">Machine Learning</span>
    <span class="cat-desc">Descripcion de lo que contiene esta categoria</span>
    <span class="cat-count">0</span>   <!-- Se actualiza al agregar codigos -->
</a>
\end{lstlisting}

\subsection{Paso 4: Definir el prefijo de la nueva categoría}

Elige un prefijo de 2--4 letras para los archivos de código. Ejemplo:
\texttt{ml} para Machine Learning → los archivos serán \archivo{ml-001.html},
\archivo{ml-002.html}, etc.


% ═════════════════════════════════════════════════════════════════════════════
\section{Guía de Estilos y CSS}
% ═════════════════════════════════════════════════════════════════════════════

Todo el diseño visual está controlado por un único archivo: \archivo{styles.css}.
\textbf{Nunca uses estilos inline en los archivos HTML} (excepto en el sidebar
donde ya existen algunos). Cualquier cambio visual debe hacerse en este archivo.

\subsection{Variables de color (tema)}

Al inicio de \archivo{styles.css} se definen las variables del tema en
la pseudo-clase \texttt{:root}. Cambiar estos valores modifica el color
en todo el sitio simultáneamente:

\begin{lstlisting}[style=cssInline, caption={Variables de color en styles.css}]
:root {
    --color-fondo:          #0f1117;  /* Fondo general oscuro */
    --color-fondo-tarjeta:  #1a1d27;  /* Fondo de tarjetas */
    --color-fondo-codigo:   #12141c;  /* Fondo de bloques de código */
    --color-borde:          #2a2d3a;  /* Color de los bordes */
    --color-texto:          #e0e0e0;  /* Texto principal */
    --color-texto-claro:    #8b8fa3;  /* Texto secundario/gris */
    --color-acento:         #6c63ff;  /* Morado principal */
    --color-acento-hover:   #5a52d5;  /* Morado al hover */
    --color-verde:          #4caf50;  /* Verde (etiqueta de lenguaje) */
    --color-azul:           #2196f3;  /* Azul (links) */
    --color-naranja:        #ff9800;  /* Naranja (avisos) */
}
\end{lstlisting}

\subsection{Cambiar el color principal del tema}

Para cambiar el color morado/acento a otro color (ej: azul), cambia únicamente:

\begin{lstlisting}[style=cssInline]
--color-acento:        #2196F3;  /* Nuevo color primario */
--color-acento-hover:  #1976D2;  /* Version oscura del nuevo color */
\end{lstlisting}

\subsection{Clases CSS importantes}

\begin{longtable}{lp{9cm}}
\toprule
\textbf{Clase} & \textbf{Descripción} \\
\midrule
\endhead
\clase{sidebar}              & Barra lateral fija de 240px \\
\clase{zona-principal}       & Todo el contenido a la derecha del sidebar \\
\clase{header}               & Barra superior con título \\
\clase{contenido}            & Contenedor central (max-width: 1000px) \\
\clase{tarjeta}              & Caja de contenido principal de un código \\
\clase{tarjeta-titulo}       & Título grande del código \\
\clase{tarjeta-descripcion}  & Descripción en gris \\
\clase{etiquetas}            & Contenedor de los tags \\
\clase{etiqueta}             & Tag individual (ej: \texttt{\#etl}) \\
\clase{etiqueta-lenguaje}    & Badge verde del lenguaje (ej: python) \\
\clase{info-meta}            & Fila de autor y fecha \\
\clase{instrucciones}        & Sección de instrucciones \\
\clase{bloque-codigo}        & Caja oscura con el código \\
\clase{barra-archivo}        & Barra superior del bloque con nombre y botón copiar \\
\clase{seccion-video}        & Sección del video \\
\clase{video-wrapper}        & Contenedor responsive 16:9 \\
\clase{video-placeholder}    & Mensaje cuando no hay video \\
\clase{footer}               & Pie de página \\
\clase{grid-proyectos}       & Grid 2 columnas en páginas de categoría \\
\clase{proyecto-card}        & Tarjeta clicable de proyecto \\
\clase{badge-tipo}           & Base del badge de tipo \\
\clase{badge-notebook}       & Badge azul ``Notebook'' \\
\clase{badge-script}         & Badge naranja ``Script'' \\
\clase{badge-proyecto}       & Badge verde ``Proyecto'' \\
\clase{stats-grid}           & Grid de 3 estadísticas en el dashboard \\
\clase{dash-categorias}      & Grid 2×2 de categorías en el dashboard \\
\clase{dash-recientes}       & Lista de códigos recientes en el dashboard \\
\clase{sidebar-categoria}    & Link de categoría en el sidebar \\
\clase{sidebar-categoria activa} & Link resaltado (categoría activa) \\
\bottomrule
\end{longtable}

\subsection{Tipos de proyecto y sus badges}

Existen exactamente tres tipos de proyecto. Usa el que corresponda:

\begin{center}
\begin{tabular}{lll}
\toprule
\textbf{Tipo} & \textbf{Clase CSS} & \textbf{Color} \\
\midrule
Proyecto  & \clase{badge-proyecto}  & Verde \\
Notebook  & \clase{badge-notebook}  & Azul \\
Script    & \clase{badge-script}    & Naranja \\
\bottomrule
\end{tabular}
\end{center}

En la página de detalle del código, el tipo se pone como texto simple
(sin clase de badge):

\begin{lstlisting}[style=htmlstyle]
<span class="etiqueta-lenguaje">python</span>
<span style="color: #8b8fa3; font-size: 14px;">Notebook</span>
\end{lstlisting}

En la página de categoría, sí usa la clase de badge:

\begin{lstlisting}[style=htmlstyle]
<span class="badge-tipo badge-notebook">Notebook</span>
\end{lstlisting}


% ═════════════════════════════════════════════════════════════════════════════
\section{Modificar el Dashboard (index.html)}
% ═════════════════════════════════════════════════════════════════════════════

\subsection{Actualizar las estadísticas}

Las tarjetas de estadísticas en el dashboard muestran 3 números:

\begin{lstlisting}[style=htmlstyle]
<div class="stats-grid">
    <div class="stat-card">
        <div class="stat-icono morado"></div>
        <div>
            <div class="stat-numero">16</div>         <!-- Total de codigos -->
            <div class="stat-label">Códigos Totales</div>
        </div>
    </div>
    <div class="stat-card">
        ...
        <div class="stat-numero">4</div>              <!-- Total de categorias -->
        ...
    </div>
    <div class="stat-card">
        ...
        <div class="stat-numero">2</div>              <!-- Total de lenguajes -->
        ...
    </div>
</div>
\end{lstlisting}

\subsection{Cambiar el nombre del repositorio}

El título que aparece en la barra superior del dashboard:

\begin{lstlisting}[style=htmlstyle]
<div class="header">
    <span class="logo">Repositorio de Demos</span>      <!-- Cambiar aqui -->
    <span class="descripcion">Repositorio Interno de Código</span>
</div>
\end{lstlisting}

\subsection{Actualizar el conteo de las categorías}

El número circular en cada tarjeta de categoría:

\begin{lstlisting}[style=htmlstyle]
<span class="cat-count">4</span>    <!-- Numero de codigos en esa categoria -->
\end{lstlisting}


% ═════════════════════════════════════════════════════════════════════════════
\section{Eliminar un Código}
% ═════════════════════════════════════════════════════════════════════════════

Para eliminar un código existente sin dejar referencias rotas:

\begin{enumerate}[leftmargin=1.5em]
    \item Elimina el archivo \archivo{codigos/xxx-00N.html}.
    \item En la página de categoría correspondiente (\archivo{categorias/xxx.html}),
          elimina la tarjeta \clase{proyecto-card} que apuntaba a ese archivo.
    \item En \archivo{index.html}, si el código aparecía en \clase{dash-recientes},
          elimina también ese \etiquetaHtml{li}.
    \item Actualiza todos los contadores mencionados en la sección 6.8.
\end{enumerate}

\begin{cajaAviso}
    Antes de eliminar, verifica con \texttt{Ctrl+F} en todos los archivos HTML
    si el nombre del archivo aparece referenciado en algún otro lugar.
\end{cajaAviso}


% ═════════════════════════════════════════════════════════════════════════════
\section{Buenas Prácticas de Edición}
% ═════════════════════════════════════════════════════════════════════════════

\subsection{Reglas generales}

\begin{itemize}[leftmargin=1.5em]
    \item \textbf{Nunca edites \archivo{styles.css} para un caso específico.}
          Los estilos son globales. Si necesitas algo puntual, usa un atributo
          \texttt{style=""} solo como último recurso.
    \item \textbf{No cambies la estructura del sidebar} en archivos individuales.
          El sidebar es idéntico en todas las páginas (excepto la clase \texttt{activa}).
    \item \textbf{Usa siempre rutas relativas.} Nunca rutas absolutas.
    \item \textbf{No incluyas librerías externas} (CDN, npm, etc.) sin aprobación
          del equipo. El sitio debe funcionar sin internet.
    \item \textbf{Prueba siempre} abriendo el archivo en un navegador antes de
          considerar el trabajo terminado.
\end{itemize}

\subsection{Consistencia de metadatos}

\begin{itemize}[leftmargin=1.5em]
    \item \textbf{Fechas:} Siempre en formato \texttt{YYYY-MM-DD} (ej: \texttt{2025-06-01}).
    \item \textbf{Tags:} En minúsculas con guiones (kebab-case): \texttt{\#fraud-detection}.
    \item \textbf{Autores:} Usa el nombre del equipo, no personas individuales
          (ej: ``ML Engineering'', no ``Juan Pérez'').
    \item \textbf{Títulos:} En español con mayúscula inicial. Sin punto final.
    \item \textbf{Descripciones cortas:} 1--2 oraciones máximo. No repitas el título.
    \item \textbf{Instrucciones:} Sé específico: qué formato de datos de entrada,
          qué retorna, qué librerías necesita.
\end{itemize}

\subsection{Editores recomendados}

\begin{itemize}[leftmargin=1.5em]
    \item \textbf{VS Code} con extensiones: \textit{Prettier}, \textit{HTML CSS Support}.
    \item \textbf{Sublime Text} con paquete \textit{Emmet}.
    \item Cualquier editor con soporte de HTML y vista previa (o usa la extensión
          \textit{Live Server} en VS Code para previsualizar en tiempo real).
\end{itemize}


% ═════════════════════════════════════════════════════════════════════════════
\section{Problemas Comunes y Soluciones}
% ═════════════════════════════════════════════════════════════════════════════

\begin{longtable}{p{5.5cm}p{8.5cm}}
\toprule
\textbf{Problema} & \textbf{Causa / Solución} \\
\midrule
\endhead

El CSS no carga (página sin estilos)
&
La ruta en \texttt{<link rel="stylesheet">} es incorrecta.
En \archivo{codigos/}, debe ser \texttt{href="../styles.css"}.
En \archivo{categorias/}, también \texttt{href="../styles.css"}.
En \archivo{index.html}, solo \texttt{href="styles.css"}.
\\[8pt]

Los links del sidebar llevan a páginas en blanco
&
La ruta en los \texttt{href} es incorrecta. Desde \archivo{codigos/},
las categorías se acceden con \texttt{../categorias/analytics.html}.
\\[8pt]

El código aparece como texto plano sin colores
&
Es comportamiento esperado si no se agregaron los \texttt{<span>}.
Para colorear, sigue la guía de la sección 5.
\\[8pt]

El botón ``Copiar'' no funciona
&
El sitio debe abrirse desde un servidor (o \texttt{localhost}), no
directamente desde el sistema de archivos (\texttt{file://}) en algunos navegadores.
Usa la extensión \textit{Live Server} de VS Code, o instala Python y ejecuta
\texttt{python -m http.server 8000}.
\\[8pt]

El video de YouTube no carga
&
Verifica que usas la URL de \textbf{embed} (\texttt{youtube.com/embed/ID})
y no la URL normal (\texttt{youtube.com/watch?v=ID}).
\\[8pt]

Aparecen caracteres como ``<'' o ``>'' rotos en el código
&
Son caracteres especiales HTML. Usa \texttt{\&lt;} para \texttt{<}
y \texttt{\&gt;} para \texttt{>}. Ver sección 5.4.
\\[8pt]

La página de categoría no muestra el nuevo código
&
No se agregó la tarjeta en el \clase{grid-proyectos} del archivo
de categoría correspondiente. Ver sección 6.7.
\\[8pt]

El número de códigos en el sidebar es incorrecto
&
El texto ``N códigos disponibles'' es estático y debe actualizarse
manualmente en todos los archivos. Ver sección 6.8.
\\[8pt]

El video local no reproduce
&
Verifica que el archivo \texttt{.mp4} existe en \archivo{videos/}
y que la ruta en \texttt{<source src="">} es correcta (\texttt{../videos/mi\_video.mp4}).
\\

\bottomrule
\end{longtable}


% ═════════════════════════════════════════════════════════════════════════════
\section{Servidor Local de Desarrollo}
% ═════════════════════════════════════════════════════════════════════════════

Aunque el sitio es estático, algunas funciones (como el botón Copiar)
requieren que se sirva a través de un servidor HTTP en lugar de abrir los
archivos directamente con el explorador de archivos.

\subsection{Opción A: Python (recomendado, sin instalación extra)}

Si tienes Python instalado, ejecuta desde la raíz del proyecto:

\begin{cajaCodigo}
\# Python 3:\\
\texttt{python -m http.server 8000}\\[4pt]
\# Python 2:\\
\texttt{python -m SimpleHTTPServer 8000}
\end{cajaCodigo}

Luego abre el navegador en: \texttt{http://localhost:8000}

\subsection{Opción B: VS Code Live Server}

\begin{enumerate}[leftmargin=1.5em]
    \item Instala la extensión \textit{Live Server} de Ritwick Dey.
    \item Abre la carpeta del proyecto en VS Code.
    \item Haz clic derecho en \archivo{index.html} → \textit{Open with Live Server}.
    \item El sitio abre automáticamente en \texttt{http://127.0.0.1:5500}.
\end{enumerate}

\subsection{Opción C: Node.js}

\begin{cajaCodigo}
\texttt{npx serve .}
\end{cajaCodigo}


% ═════════════════════════════════════════════════════════════════════════════
\section{Lista de Verificación para Nuevas Entradas}
% ═════════════════════════════════════════════════════════════════════════════

Usa esta lista antes de considerar completo cualquier nuevo código:

\textbf{Archivo de código (\archivo{codigos/xxx-00N.html}):}
\begin{itemize}[leftmargin=1.5em, label={\(\square\)}]
    \item El \texttt{<title>} tiene el nombre correcto del proyecto.
    \item El sidebar tiene la categoría correcta marcada como \texttt{activa}.
    \item El \clase{tarjeta-titulo} está completo.
    \item El lenguaje en \clase{etiqueta-lenguaje} es correcto.
    \item El tipo (Proyecto / Notebook / Script) está correcto.
    \item La descripción en \clase{tarjeta-descripcion} es clara y concisa.
    \item Hay al menos 2 etiquetas en \clase{etiquetas}.
    \item El autor y la fecha en \clase{info-meta} son correctos (\texttt{YYYY-MM-DD}).
    \item Las instrucciones de uso explican entradas, salidas y dependencias.
    \item El nombre del archivo en \clase{barra-archivo} es el correcto (\texttt{.py}).
    \item El código es el real (no el placeholder del motor de recomendaciones).
    \item Si hay video: la URL de YouTube o la ruta del \texttt{.mp4} son correctas.
    \item El botón ``Copiar'' funciona al abrir en el navegador.
\end{itemize}

\textbf{Página de categoría:}
\begin{itemize}[leftmargin=1.5em, label={\(\square\)}]
    \item Se agregó la tarjeta en \clase{grid-proyectos}.
    \item El \texttt{href} de la tarjeta apunta al archivo correcto.
    \item El badge de tipo (\texttt{badge-notebook}, etc.) es correcto.
    \item El contador \clase{categoria-contador} se actualizó.
\end{itemize}

\textbf{Dashboard (\archivo{index.html}):}
\begin{itemize}[leftmargin=1.5em, label={\(\square\)}]
    \item El número total en \clase{stat-numero} se actualizó.
    \item El \clase{cat-count} de la categoría se actualizó.
    \item El texto ``N códigos disponibles'' en el sidebar footer se actualizó.
\end{itemize}

\textbf{Todos los archivos:}
\begin{itemize}[leftmargin=1.5em, label={\(\square\)}]
    \item El texto ``N códigos disponibles'' en el \clase{sidebar-footer}
          de todos los archivos de \archivo{codigos/} se actualizó.
\end{itemize}


% ─── FIN DEL DOCUMENTO ───────────────────────────────────────────────────────
\vspace{2cm}
\noindent\textcolor{cpAccent}{\rule{\linewidth}{0.6pt}}

\vspace{0.5cm}
\begin{center}
    {\small\color{cpTextoClaro}
    ConversionPulse · Repositorio Interno de Código · 2025 \\[4pt]
    Documentación generada para uso interno del equipo.}
\end{center}

\end{document}
